\begin{problem}[Presheaves on the Sierpiński space]
Let $X=\{\eta,s\}$ with opens $\varnothing,\{\eta\},X$.  Classify normalized presheaves of abelian groups on $X$ up to isomorphism.
\end{problem}

\paragraph{Why this problem matters.}
This is the smallest nontrivial topological example and turns presheaves into explicit algebra.

\paragraph{Useful ideas before starting.}
Normalize $\mathcal F(\varnothing)=0$ and list the remaining inclusions.

\paragraph{Guided hints.}
Normalize $\mathcal F(\varnothing)=0$ and list the remaining inclusions.

\begin{solution}
A presheaf is exactly two abelian groups $A=\mathcal F(X)$ and $B=\mathcal F(\{\eta\})$ together with one homomorphism $r:A\to B$.  There is no further independent restriction, since maps to $0$ are unique.  An isomorphism of presheaves is exactly a commutative square with vertical group isomorphisms.  Hence isomorphism classes are isomorphism classes of arrows in $\mathbf{Ab}$.
\end{solution}

\paragraph{What happened mathematically.}
A finite topology converts the abstract presheaf definition into a diagram category.

\paragraph{Extensions.}
Classify sheaves on the same space in VI/13 and compare.

\paragraph{Provenance.}
New canonical problem, motivated by the finite-space viewpoint.
